\documentclass[epsfig,a4paper,11pt,titlepage,twoside,openany]{book}
\usepackage{epsfig}
\usepackage{plain}
\usepackage{setspace}
\usepackage[paperheight=29.7cm,paperwidth=21cm,outer=1.5cm,inner=2.5cm,top=2cm,bottom=2cm]{geometry} % per definizione layout
\usepackage{titlesec} % per formato custom dei titoli dei capitoli
\usepackage[]{todonotes}
\usepackage[]{lineno}
\usepackage{hyperref}
\usepackage{cleveref}
\usepackage{enumitem}
\usepackage{tabularray}
\usepackage{csvsimple}
\usepackage{adjustbox}
\usepackage{pgfplotstable}
\usepackage{booktabs}
\usepackage{xcolor}
\usepackage{colortbl}
\usepackage{array}
\usepackage{float}
\usepackage{tabularx}

\setlist{noitemsep,topsep=0pt}
\linenumbers


\usepackage[utf8]{inputenc} % per UTF-8.

\singlespacing

\usepackage[english]{babel}

\hyphenation{WhatsApp}



% CMDs
\newcommand{\tdin}[1]{\todo[color=blue!30,inline]{TODO: #1}}

\begin{document}

% nessuna numerazione
\pagenumbering{gobble}
\pagestyle{plain}

\thispagestyle{empty}

\begin{center}
    \begin{figure}[h!]
        \centerline{\psfig{file=marchio_unitrento_colore_it_202002.eps,width=0.6\textwidth}}
    \end{figure}

    \vspace{2 cm}

    \LARGE{Department of Information Engineering and Computer Science\\}

    \vspace{1 cm}
    \Large{Bachelor Degree in\\
        Computer Science
    }

    \vspace{2 cm}
    \Large\textsc{Final Dissertation\\}
    \vspace{1 cm}
    \Huge\textsc{Autonomous Performance Benchmarking of Large Language Models for Online Retail Support Systems\\}
    % \Large{\it{Sottotitolo (alcune volte lungo - opzionale)}}

    \vspace{2 cm}
    \begin{tabular*}{\textwidth}{ c @{\extracolsep{\fill}} c }
        \Large{Supervisor} & \Large{Student}\\
        \Large{Marco Patrignani}& \Large{Giovanni Feltrin}\\
    \end{tabular*}

    \vspace{2 cm}

    \Large{Academic year 2024/2025}

\end{center}


\clearpage

\thispagestyle{empty}

\begin{center}
    {\bf \Huge Acknowledgements}
\end{center}

\vspace{4cm}

\emph{
    ...thanks to
}

\clearpage
\pagestyle{plain} % nessuna intestazione e pie pagina con numero al centro

% inizio numerazione pagine in numeri arabi
\mainmatter

% indice
\tableofcontents
\clearpage

% gruppo per definizone di successione capitoli senza interruzione di pagina
\begingroup
% nessuna interruzione di pagina tra capitoli
% ridefinizione dei comandi di clear page
\renewcommand{\cleardoublepage}{}
\renewcommand{\clearpage}{}
% redefinizione del formato del titolo del capitolo
% da formato
%   Capitolo X
%   Titolo capitolo
% a formato
%   X   Titolo capitolo

\titleformat{\chapter}
{\normalfont\Huge\bfseries}{\thechapter}{1em}{}

\titlespacing*{\chapter}{0pt}{0.59in}{0.02in}
\titlespacing*{\section}{0pt}{0.20in}{0.02in}
\titlespacing*{\subsection}{0pt}{0.10in}{0.02in}

\chapter{Introduction}
\label{cha:introduction}

% \section{Abstract}

% Context

E-commerce businesses rely heavily on email marketing to drive revenue growth, but emails have become increasingly saturated and suffer from declining engagement rates.
This has created a pressing need for these businesses to find more effective communication channels to reach their customers.
The messaging platform \textit{WhatsApp} recently released APIs to interact with the platform in a systematic manner, notably the ability to automatically send messages.
This enables e-commerce businesses to leverage WhatsApp as an alternative communication channel that offers significantly higher open and response rates than email while supporting the same marketing capabilities including automated campaigns, customer segmentation, and targeted messaging.
E-commerce businesses can easily establish automated messaging workflows, and customers frequently engage with these communications by responding to the messages.

% Problem

However, this increased engagement presents both opportunities and challenges: e-commerce businesses must now manage a substantially larger volume of customer inquiries, many of which are tedious and do not require specialized knowledge nor complex actions.
WhatsApp has platform-specific limitations that worsen this challenge: it enforces strict restrictions on message responses, particularly limiting businesses from responding to customer messages after a 24-hour window has elapsed since the last customer-initiated message \cite{wa_policy}.

% Solution

To address these issues, this thesis presents the implementation of an LLM-powered assistant for e-commerce customer service.
The assistant autonomously guides customers through the purchase process by suggesting relevant products, answering order inquiries about contents and shipping status, and providing business-specific information.
The system also aids human intervention by prioritizing conversations and categorizing conversations to enable different departments to operate independently and efficiently.
The architecture we chose for the assistant employs a 2-phase approach comprising \textit{Context Gathering}, the first, and \textit{Answer Generation}, the second.
This architectural necessity enables the system to efficiently handle large product catalogs and order histories by first identifying relevant information and then generating appropriate responses based on that context.
Since implementing such a system can be accomplished through various platforms, including no-code solutions, the primary contribution of this thesis lies in the validation.

% Validation

In particular, we focus on developing a systematic process that can automatically and objectively evaluate its performance, in function of the models chosen to complete the tasks.
To do this, we test each model using a set of synthetic chat conversations and evaluate the results using another LLM that acts as a judge (technique known as LLM-as-a-Judge).
We carefully engineered the prompt used for the judge model to match human evaluation standards, enabling consistent and reliable assessment of model performance across all test scenarios.
We validate each of the two phases separately, leveraging their independent operational capability (by providing standardized, synthetically generated context to the \textit{Answer Generation} phase).
Evaluation covers different dimensions for each phase: for the Context Gathering phase, we evaluate the ability of the model to correctly choose what context is needed and how to obtain it; for the Answer Generation Phase, we evaluate the actual message response together with chat categorization and prioritization.
Additionally, both phases are evaluated across three common performance metrics: failure rate, cost, and speed.

Our results show that for the Context Gathering phase, which is a simpler task, smaller models performed best in terms of speed and cost while still achieving very good accuracy.
For the Answer Generation phase, however, smaller models struggled to maintain quality and produced significantly more hallucinations.
Among the larger models, there is no single best choice: organizations must balance between achieving the highest accuracy, minimizing costs, and optimizing response speed based on their specific requirements.

\section{Outline}

We begin by establishing the theoretical foundations in \cref{cha:background}, covering the key technologies and concepts that underpin our assistant implementation.

\Cref{cha:impl} presents the core implementation of the assistant, detailing our 2-phase architecture that addresses computational cost scalability challenges.
The Context Gathering phase (\cref{sec:context-gathering}) analyzes conversation history to determine what contextual information is needed, while the Answer Generation phase (\cref{sec:answer-generation}) uses this context to generate priority levels, mood classification, conversation categories, and actions.
We also discuss output format considerations (\cref{sec:output_format}) across different LLM providers and reliability mechanisms (\cref{sec:reliability}) including retry logic for handling model inconsistencies.

\Cref{cha:validation} presents our systematic evaluation framework using synthetic chat conversations to assess multiple LLM models across major providers.
We employ an LLM-as-a-Judge methodology (\cref{sec:llm-judge}) for automated evaluation, testing both phases independently to measure accuracy, performance, and cost-effectiveness.
The results (\cref{sec:cg-results}, \cref{sec:ag-results}) show that the first phase can be handled by most models, while the more complex second phase requires careful consideration of trade-offs between quality, speed, and cost.

\Cref{cha:related-work} examines how this thesis builds upon and extends existing research in RAG techniques and LLM-as-a-Judge methodologies, with particular focus on their application to customer service automation.

Finally, \cref{cha:conclusion} presents our considerations and potential extensions to enhance the capabilities of the system.


\newpage
\chapter{Background}
\label{cha:background}

This chapter establishes the theoretical foundations necessary for understanding the implementation and evaluation of an LLM-powered customer service assistant for e-commerce applications. We begin by examining the capabilities and limitations of Large Language Models in practical applications, followed by an analysis of the specific contextual requirements for e-commerce customer service automation. We then explore Retrieval Augmented Generation (RAG) techniques that enable LLMs to access external knowledge sources, and conclude with an overview of vector database technologies that provide the semantic search capabilities essential for efficient information retrieval in our system architecture.
\tdin{Check updates}

\section{Large Language Models}

Large Language Models (LLMs) are neural network architectures trained to understand and generate natural language responses. These models demonstrate strong performance across general knowledge tasks but require domain-specific context to be effectively applied in specialized applications such as e-commerce customer service.

The fundamental limitation of LLMs for practical applications lies in their training data cutoff and lack of access to real-time, business-specific information.
While these models possess broad linguistic capabilities and general world knowledge, they cannot access current product catalogs, order statuses, or company-specific policies without external information retrieval mechanisms.

\section{Domain-Specific Context in the E-commerce Niche}

This thesis focuses on automating customer service responses that primarily involve retrieving and presenting existing information rather than executing transactional actions.
We designed the automated assistant to handle informational queries by processing and presenting data already stored within the e-commerce system, emphasizing information delivery over system modifications or direct transaction processing.
In particular, three main categories of information are required to effectively respond to customer inquiries:
\begin{itemize}
    \item Order history information for questions regarding order status and contents
    \item Available product information to guide customers through the purchase process and provide relevant product recommendations
    \item E-commerce business rules including accepted payment methods, available delivery options, and courier information
\end{itemize}

\section{Retrieval Augmented Generation}
\label{sec:rag}

RAG is a technique that combines information retrieval with language generation to provide more accurate and contextual responses.
The implementation of RAG follows a straightforward process: when a user submits a query, the system first converts the query into a vector representation and searches a vector database for semantically similar content.
The most relevant information is then retrieved and provided as context to the LLM, which generates a response that incorporates this retrieved knowledge alongside its pre-trained capabilities.
This approach allows LLMs to access up-to-date, domain-specific information that wasn't present in their original training data.

\section{Vector Databases}
\label{sec:vecdb}

Vector databases are specialized storage systems designed to efficiently store and retrieve multidimensional vector representations of data; in our case, for example, text embeddings generated by machine learning models.
They enable semantic search capabilities by measuring similarity between vectors, allowing systems to find contextually relevant information even when exact keyword matches don't exist.

\section{Prompt Engineering}
\label{sec:prompt-engineering}

Prompt engineering is the practice of designing and optimizing input prompts to guide LLM behavior and improve response quality.
Effective prompts provide clear instructions, appropriate context, and formatting guidelines that help the model generate more accurate and relevant responses.
In customer service applications, well-crafted prompts can significantly improve the consistency and helpfulness of automated responses by establishing the assistant's role, communication style, and response structure.

\section{Chain-of-Thought and Reasoning Capabilities}
\label{sec:thinking}

Chain-of-Thought (CoT) prompting is a technique that encourages LLMs to break down complex problems into sequential reasoning steps, making their decision-making process more transparent and often more accurate.
By explicitly requesting the model to "think through" a problem step-by-step, CoT prompting can improve performance on tasks requiring logical reasoning, multi-step analysis, or complex decision-making.
Modern LLMs increasingly incorporate built-in reasoning capabilities, often referred to as "thinking" modes, where the model internally processes information through multiple reasoning steps before generating a final response.
These capabilities allow models to perform more thorough analysis of queries, consider multiple aspects of a problem, and provide more thoughtful responses.
The trade-off for enhanced reasoning capabilities typically involves increased computational cost and longer response times, as models spend additional processing cycles on internal reasoning before producing output.


\newpage
\chapter{Implementation}
\label{cha:impl}

This chapter presents the core implementation of our LLM-powered assistant for e-commerce customer service, detailing the 2-phase architecture designed to address computational cost scalability challenges while maintaining high response quality.
We examine the Context Gathering and Answer Generation phases, discuss output format considerations across different LLM providers, and describe the reliability mechanisms implemented to handle model inconsistencies.

\section{The 2-Phase Architecture}

An analysis of the pricing structures for Large Language Models used in this thesis \cite{gemini_pricing, anthropic_pricing, openai_pricing} demonstrates that computational costs increase proportionally with the volume of context tokens processed by the model.
This observation highlighted a fundamental scalability challenge: most e-commerce platforms maintain catalogs containing hundreds to thousands of products and variations, making it impractical to include all available information in every model inference.

To address this limitation we implemented a 2-phase architecture, where the Context Gathering Phase receives the conversation history as input and decides which contextual elements should be included in the subsequent Answer Generation Phase.
This approach represents a more sophisticated form of RAG (\cref{sec:rag}) that ensures older or resolved conversations do not degrade output quality by selectively retrieving only the most relevant contextual information.
This approach offers significant scalability advantages and provides a framework that can be extended to incorporate additional context categories as system requirements evolve.

\section{Output Format Considerations}
\label{sec:output_format}

A critical design consideration for the implementation involves the output structure format. Since the LLM's output undergoes automatic parsing, it must maintain consistent structure with predefined fields that cannot be left to the model's discretion.

Various LLM providers offer different levels of output formatting control. Some models provide a \textit{structured output} mode, which accepts a JSON schema and constrains the model's response to conform to the specified structure. Other models offer only a \textit{JSON mode} that enforces JSON formatting without a predefined schema structure.
Some models do not support structured output formatting at all, requiring reliance solely on prompt engineering techniques.

For this implementation, JSON was selected as the output format over alternatives such as XML due to its widespread support across different model providers.
Models that do not support JSON mode (notably those from Anthropic) or those with intermediate capabilities (such as older OpenAI models) require additional prompt engineering.
These models receive a TypeScript representation of the desired output format within their prompt, which provides a more compact specification than JSON Schema while maintaining sufficient flexibility for the system's requirements.
Since these models are trained to output code formatted as Markdown, they typically wrap JSON responses in code blocks, requiring the removal of Markdown formatting from the beginning and end of the response.
While prompt engineering techniques were explored to eliminate this behavior during development, they proved inconsistent across different queries and contexts, whereas simple string trimming provides a reliable solution for extracting the JSON content.

\section{System Architecture Details}

\subsection{Phase 1}

The Context Gathering Phase utilizes a language model that has been specifically configured to analyze conversational context and determine the relevant information required for generating an appropriate response.
The output structure of this phase consists of two primary components: a \textit{boolean} value that indicates whether order history should be incorporated into the context, and a nullable \textit{string} parameter that contains a search query for product retrieval.
When generating such queries, the language model is instructed to formulate queries optimized for vector database similarity search operations (\cref{sec:vecdb}).

This phase could also be used to synthesize the chat history to optimize Phase 2 by providing a condensed summary of the conversation; however, in this thesis we focus on shorter conversations and maximum accuracy, prioritizing the preservation of complete conversational context over computational efficiency gains from summarization.

\subsection{Intermediate Phase}

When a query is generated during Phase 1, the same embedding model that was employed for the product catalog indexing converts it into an embedding, ensuring consistency in the vector space representation.
A similarity search is then performed against the product catalog embeddings, with the top-N most relevant products being selected based on cosine similarity scores (a threshold is also used to avoid including unrelated information).
These selected products, along with order information when requested, are subsequently formatted and passed to the Answer Generation Phase.

\subsection{Phase 2}

The Answer Generation Phase receives the contextual information assembled from the previous phases along with the complete conversation history.
In this phase, \cref{sec:prompt-engineering} has a much greater impact on the quality and effectiveness of the generated responses than in Phase 1.
The output structure consists of multiple components:
\begin{itemize}
    \item A \textit{Priority} enumeration that indicates the urgency level of the conversation (ranging from no urgency to urgent human intervention required)
    \item A \textit{Mood} enumeration that captures the customer's emotional state
    \item A list of \textit{categories} that tag the conversation with relevant business classifications
    \item An \textit{actions} array (that might be empty) containing the actions the model choose to take and their parameters
\end{itemize}
The use of enumerations for \textit{Priority} and \textit{Mood} provides significant advantages as they effectively constrain the model's output to predefined, valid options, ensuring consistency and reliability in the structured responses while preventing the generation of unexpected or malformed values.

The actions array represents a flexible and scalable architectural design that enables future system expansion.
Currently, the system implements a single action type for sending messages; however, the framework can be trivially extended to incorporate additional action types such as tag modification or checkout creation without requiring fundamental architectural changes.
This structured approach also enhances the model's awareness of its operational context and decision-making process.
During the development phase, empirical observations revealed that models would sometimes misuse an old \textit{message} field for commentary or justification purposes rather than customer-facing communication, a behavior that the structured action framework effectively eliminates.

\section{Model Reliability and Error Mitigation}

Both model calls in the two-phase architecture implement a retry mechanism with a maximum limit of three attempts to mitigate potential failures.
The vast majority of failures encountered during system operation occur during the deserialization process, specifically in instances where language models fail to adhere to the specified output structure formats.

To ensure fair evaluation across different models, the cumulative costs and execution times across multiple retry attempts are recorded and summed for each query.
This approach provides an accurate representation of the total computational overhead required by models that demonstrate varying levels of consistency in output format adherence.

Empirical testing conducted during the development phase revealed that flagship models achieved error rates of less than 1\% in output format compliance.
Models equipped with structured output capabilities demonstrated effectively zero error rates, highlighting the significant reliability improvements offered by constrained generation features.
While this implementation detail will not be examined in depth in the validation chapter, its effects on system performance are reflected in the measured costs and execution times presented in the results.


\newpage
\chapter{Validation}
\label{cha:validation}

\begin{itemize}
    \item Output Format
          \begin{itemize}
              \item JSON and non-JSON models
              \item Retries, accuracy
          \end{itemize}
    \item Possible Actions
          \begin{itemize}
              \item SendMessage
          \end{itemize}
    \item Categories
    \item Priority
          \begin{itemize}
              \item Should account for its answer
          \end{itemize}
    \item Evaluation
          \begin{itemize}
              \item Automated evaluation of multiple models
              \item Includes price, speed, accuracy
          \end{itemize}
\end{itemize}


\newpage
\chapter{Related Work}
\label{cha:related-work}

This thesis presents a specialization of existing techniques, along with targeted improvements for our specific use case.
For example, our Context Gathering phase, as it represents a specialized implementation of RAG (Retrieval Augmented Generation).
As Y.Gao et al. \cite{gao2023retrieval} demonstrate, numerous RAG techniques exist; however, our version is optimized for continuous support chat environments by filtering out outdated, resolved, and irrelevant data to focus exclusively on the customer's current needs.
Another example is our application of the \textit{LLM-as-a-Judge} methodology: D. Li et al. \cite{llmasajudge_applications} provide a comprehensive taxonomy of research in this area.
This thesis contributes a practical implementation with evaluation results, featuring specialized evaluation criteria tailored specifically to the e-commerce domain.

In examining similar approaches to LLM deployment in customer service, Wulf et al. \cite{wulf2024exploringpotentiallargelanguage} present a practical implementation for telecommunications support, demonstrating feasibility through prototyping with actual customer data.
Their cognitive task categorization parallels our Context Gathering and Answer Generation phases.
While they focus on technical feasibility and manual validation, our implementation emphasizes autonomous evaluation via LLM-as-a-Judge methodology and business metrics, extending to comprehensive cost analysis and deployment scenarios for e-commerce.
This addresses their identified gap for large-scale validation by providing a framework for model selection based on operational constraints and profit margins.


\newpage
\chapter{Conclusion}
\label{cha:conclusion}


\endgroup

% bibliografia in formato bibtex
%
% aggiunta del capitolo nell'indice
\addcontentsline{toc}{chapter}{Bibliography}
% stile con ordinamento alfabetico in funzione degli autori
\bibliographystyle{plain}
\bibliography{biblio}
%%%%%%%%%%%%%%%%%%%%%%%%%%%%%%%%%%%%%%%%%%%%%%%%%%%%%%%%%%%%%%%%%%%%%%%%%%
%%%%%%%%%%%%%%%%%%%%%%%%%%%%%%%%%%%%%%%%%%%%%%%%%%%%%%%%%%%%%%%%%%%%%%%%%%
%% Nota
%%%%%%%%%%%%%%%%%%%%%%%%%%%%%%%%%%%%%%%%%%%%%%%%%%%%%%%%%%%%%%%%%%%%%%%%%%
%% Nella bibliografia devono essere riportati tutte le fonti consultate
%% per lo svolgimento della tesi. La bibliografia deve essere redatta
%% in ordine alfabetico sul cognome del primo autore.
%%
%% La forma della citazione bibliografica va inserita secondo la fonte utilizzata:
%%
%% LIBRI
%% Cognome e iniziale del nome autore/autori, la data di edizione, titolo, casa editrice, eventuale numero dell’edizione.
%%
%% ARTICOLI DI RIVISTA
%% Cognome e iniziale del nome autore/autori, titolo articolo, titolo rivista, volume, numero, numero di pagine.
%%
%% ARTICOLI DI CONFERENZA
%% Cognome e iniziale del nome autore/autori (anno), titolo articolo, titolo conferenza, luogo della conferenza (città e paese), date della conferenza, numero di pagine.
%%
%% SITOGRAFIA
%% La sitografia contiene un elenco di indirizzi Web consultati e disposti in ordine alfabetico.
%% E’ necessario:
%%   Copiare la URL (l’indirizzo web) specifica della pagina consultata
%%   Se disponibile, indicare il cognome e nome dell’autore, il titolo ed eventuale sottotitolo del testo
%%   Se disponibile, inserire la data di ultima consultazione della risorsa (gg/mm/aaaa).
%%%%%%%%%%%%%%%%%%%%%%%%%%%%%%%%%%%%%%%%%%%%%%%%%%%%%%%%%%%%%%%%%%%%%%%%%%
%%%%%%%%%%%%%%%%%%%%%%%%%%%%%%%%%%%%%%%%%%%%%%%%%%%%%%%%%%%%%%%%%%%%%%%%%%

\titleformat{\chapter}
{\normalfont\huge\bfseries}{Attachments \thechapter}{1em}{}
\appendix
\chapter{Chat Interaction Examples}

\section{Comparison of Gemini 2.5 Pro and Gemini 2.5 Flash}
\label{sec:pro_vs_flash}

\subsection*{Messages}
\begin{itemize}
    \item \textbf{Customer:} Hello, I need to order some supplies for my office. Can you give me a bulk price quote for 50 reams of your standard printer paper and 200 of your basic black ballpoint pens?
\end{itemize}

\subsection*{Available Products}
\begin{itemize}
    \item \textbf{Standard 8.5x11 Printer Paper}
          \begin{itemize}
              \item Description: 92 brightness, 20 lb. multipurpose printer paper. Ideal for everyday printing. Sold by the ream (500 sheets).
              \item Price: 8.99€
              \item URL: [url omitted]
              \item Variations: 1 Ream - 8.99€
          \end{itemize}

    \item \textbf{Basic Black Ballpoint Pens}
          \begin{itemize}
              \item Description: Reliable and smooth-writing black ballpoint pens with a medium point. Perfect for office use. Sold in boxes of 12.
              \item Price: 5.99€
              \item URL: [url omitted]
              \item Variations: Box of 12 - 5.99€
          \end{itemize}
\end{itemize}

\subsection*{Model Responses}

\begin{table}[H]
    \begin{tabularx}{\textwidth}{|l|X|}
        \hline
        \textbf{Model}   & \textbf{Response}                                                                                                                                                                                                                                                                                                                                                                                                                                                                                                                                                        \\
        \hline
        Gemini 2.5 Pro   & Hi there! Happy to help you with that order. Here are the products you mentioned: - 50 reams of Standard 8.5x11 Printer Paper: [url omitted] - For the pens, they come in boxes of 12, so you'd need 17 boxes for 200 pens: [url omitted] \textcolor{red}{I've passed your message on to our sales team} to see what they can do for you regarding a bulk discount. \textcolor{red}{Someone will be in touch} with a quote soon!                                                                                                                                         \\
        \hline
        Gemini 2.5 Flash & Hello! I can provide you with a quote based on our current pricing for the quantities you're looking for. For 50 reams of our Standard 8.5x11 Printer Paper, the cost would be 449.50 €. You can find more details here: [url omitted]. For 200 Basic Black Ballpoint Pens, since they are sold in boxes of 12, you would need to purchase 17 boxes to get 204 pens. The total cost for 17 boxes would be 101.83 €. You can view the pens here: [url omitted]. The estimated total for your order would be 551.33 €. Please let me know if you have any other questions! \\
        \hline
    \end{tabularx}
\end{table}

\subsection*{Explanation}

The models were given specific instructions not to make claims like those highlighted in red.
This is because the AI models cannot actually pass messages to other teams or guarantee that someone will contact the customer, since it is assumed that human operators respond within the same chat interface.


\end{document}
